\chapter{Sprint 2 : Surfaces operationnelles}

\section{Introduction}
Le Sprint 1 a produit une colonne vertebrale qui fonctionnait de bout en
bout mais sans aucune surface graphique. Un proprietaire ne pouvait voir
ses donnees qu'a travers l'editeur de tables de Supabase Studio, ce qui
convient a un developpeur mais reste inutilisable pour un proprietaire de
clinique. Le Sprint 2 a donc ete consacre a la construction de chaque
surface que le proprietaire et le patient touchent au quotidien : le
tableau de bord multi-locataire avec sa barre laterale et sa barre
superieure au theme anti-gravite, le calendrier, la gestion des
rendez-vous et des clients, le pipeline specifique aux laboratoires (un
tableau Kanban avec une etape OCR d'image de prescription), et la PWA
patient avec son onboarding par OTP telephone.

\section{Backlog du sprint}
Le backlog du Sprint 2 a tire des exigences BF06--BF10 (surfaces du
tableau de bord), BF24 (OCR des prescriptions), BF29--BF30 (PWA + OTP
telephone), plus une longue liste d'exigences d'utilisabilite qui ont
emerge lors de la revue du Sprint 1.

\section{Coque du tableau de bord proprietaire}

\subsection{Langage visuel}
Le tableau de bord adopte ce que nous avons appele le theme
\emph{anti-gravite} : (a) une paire d'orbes radiaux floutes a position
fixe qui derivent lentement a travers la page (indigo et violet
principalement, avec un accent cyan), (b) un confetti de minuscules
particules qui flottent vers le haut spawnees par JS au chargement, (c)
des cartes qui se bercent doucement avec une animation
\texttt{translateY} de six secondes, (d) une pulsation lumineuse sur les
boutons d'action principale, (e) un gradient coherent
(\texttt{linear-gradient(135deg, \#6366f1, \#8b5cf6, \#22d3ee)}) utilise
pour la marque, les boutons primaires et les chiffres du minuteur de
consultation. Toutes les animations respectent
\texttt{prefers-reduced-motion}.

\subsection{Mise en page}
Le tableau de bord utilise une grille CSS a deux colonnes : une barre
laterale a largeur fixe (240 px sur ordinateur, tiroir hors-canevas en
dessous de 1024 px) et un panneau principal fluide. La barre superieure
contient le titre de la page, les boutons d'action (selecteur de theme,
bouton de scan QR, cloche, pastille utilisateur), et sur mobile un
hamburger qui bascule le tiroir.

Le contenu de la barre laterale est sensible au role. Un proprietaire
voit Vue d'ensemble, Rendez-vous, Calendrier, Clients, Equipe, Pipeline
labo (si laboratoire), Catalogue de tests (si laboratoire), Vehicules
(si garage), Services, Forfaits, Demand Fill, Schedule Insights,
Prescrire (non-labo), Anniversaires, Messages, Notifications, Test
assistant, Activite, Parametres.

\subsection{Routage}
Le tableau de bord utilise un routage cote client base sur les hashes
--- pas de \texttt{history.pushState}, pas de bibliotheque de routeur
--- implemente en environ quarante lignes de JS. Le routeur ecoute
\texttt{hashchange} et \texttt{visibilitychange}, appelle une seule
fonction \texttt{render()} qui prend \texttt{location.hash}, parse le
segment de route, et delegue a la fonction \texttt{render*()}
appropriee. Un \texttt{Promise.race} entoure l'appel d'un timeout de
huit secondes qui fait apparaitre une carte de recuperation si quelque
chose plante.

\section{Page Vue d'ensemble}
La Vue d'ensemble est la page sur laquelle le proprietaire atterrit
apres connexion. Nous l'avons reorganisee deux fois durant le projet.
La premiere version mettait la banniere IA en haut et le calendrier
pres du bas ; les tests d'utilisabilite ont revele que les
proprietaires se preoccupent surtout du \emph{jour meme} et veulent le
voir immediatement. La seconde version a place le calendrier en haut.

La disposition actuelle, du haut vers le bas : (1) pour les entreprises
non-laboratoire, une grande carte \emph{Check-in QR} avec l'action
principale \og Scanner le QR du patient \fg{} ; (2) le calendrier
mensuel complet ; (3) la banniere Receptionniste IA on/off ; (4) les
pastilles de statut (WhatsApp verifie, plan) ; (5) les tuiles KPI
(rendez-vous du jour, clients, medecins, notifications 24h) ; (6) une
section a deux colonnes : tableau des rendez-vous a venir a gauche, fil
d'activite recente a droite.

\section{Vue calendrier}
Le calendrier est une grille mensuelle (sept colonnes par cinq ou six
lignes). Chaque cellule affiche le numero du jour, la cellule du jour
courant a une bordure indigo, et jusqu'a trois pastilles de rendez-vous
listent l'heure et le nom du client. Les cellules avec plus de
rendez-vous montrent un lien \og +N de plus \fg{} qui ouvre une fenetre
modale avec la liste complete du jour.

Le menu deroulant filtre-medecin, en haut a droite, permet au
proprietaire de restreindre le calendrier a un seul praticien.

\begin{figure}[H]
  \centering
  \fbox{\parbox[c][6cm][c]{0.95\textwidth}{\centering\textit{(Inserer la capture d'ecran du calendrier)}}}
  \caption{Vue mensuelle du calendrier}
  \label{fig:calendar}
\end{figure}

\section{Gestion des rendez-vous}
La vue Rendez-vous liste chaque rendez-vous jamais reserve, ordonne
descendant par date et heure. Des filtres en haut (plage de dates,
statut, source) permettent au proprietaire de restreindre la liste.
Chaque ligne est cliquable et ouvre une fenetre modale ou le rendez-vous
peut etre annule ou annote. Deux boutons d'action rapide se trouvent
sur chaque ligne confirmee : \emph{Arrive} (horodate l'arrivee du
patient) et \emph{Termine} (cloture la visite et marque
\texttt{completed\_at}). Une fois une ligne completee, les boutons sont
remplaces par une pastille verte montrant la duree reelle de la visite
--- information que la fonctionnalite d'apprentissage de duree
exploitera plus tard.

\section{Gestion des clients}
La vue Clients est l'annuaire des patients. Chaque client qui a deja
envoye un message a l'entreprise dispose d'une ligne, avec un panneau
de detail au clic montrant son historique complet : rendez-vous,
conversations, ordres de laboratoire (si laboratoire), avis,
notifications. Le bouton Editer ouvre une fenetre modale avec nom,
telephone, age, email, anniversaire, et notes.

Un filtre de recherche en direct en haut de la liste restreint les
resultats au fur et a mesure de la frappe.

\section{Pipeline laboratoire}
Pour les locataires laboratoires, le tableau de bord montre une vue
supplementaire que la clinique n'a pas besoin : un tableau Kanban a
quatre colonnes representant les etapes de traitement des echantillons.
\textbf{Demande}, \textbf{Echantillon collecte}, \textbf{En analyse},
\textbf{Pret a envoyer}.

Chaque colonne montre une carte par lot patient (une carte peut
contenir plusieurs analyses pour la meme personne). Au sein d'une
carte, les analyses individuelles peuvent etre avancees independamment,
mais le \emph{lot} reste dans la colonne correspondant a son analyse la
plus en retard. Un bouton de progression en lot permet au technicien de
deplacer chaque patient d'une colonne a la suivante en un seul clic.

Lorsqu'un lot atteint \emph{Pret a envoyer}, une automatisation
(\texttt{lab-results-notify}) declenche un message WhatsApp au patient,
optionnellement accompagne d'un PDF detaillant les resultats.

\section{OCR des images de prescriptions}
L'etape de pre-reservation dans un laboratoire est plus exigeante que
dans une clinique. Le patient tient une prescription papier ; il a
besoin de savoir quelles analyses vont lui etre facturees, quel est le
total, si un jeune est requis. Demander au patient de taper les noms
d'analyses dans un chat est lent et source d'erreurs --- les
prescriptions manuscrites en francais utilisent souvent des
abreviations (\texttt{NFS}, \texttt{TSH}, \texttt{Glycemie a jeun},
\texttt{Bilan lipidique}) que le patient lui-meme ne sait pas lire.

La solution : lorsqu'un message entrant contient une photo et que le
locataire est un laboratoire, le webhook invoque le point de
terminaison vision de Claude Sonnet avec l'image et un prompt systeme
qui inclut le catalogue d'analyses reel du laboratoire comme
vocabulaire. Claude renvoie un JSON structure avec les noms d'analyses
reconnus, le nom du prescripteur, la date de prescription, et un score
de confiance.

Les analyses reconnues sont ensuite repassees au cerveau IA, qui les
convertit en proposition de reservation.

Une passe OCR typique prend entre quatre et huit secondes, et la
precision sur les prescriptions francaises a ete excellente : 88\%
des reconnaissances d'analyses dans notre echantillon interne ont ete
correctes au premier passage.

\section{Notifications et la cloche}
Les notifications en temps reel sont essentielles. Un proprietaire qui
fait une consultation ne peut pas constamment rafraichir le tableau de
bord ; il a besoin d'un signal passif que quelque chose s'est passe.
Nous avons superpose trois canaux de notification : (a) son (un bref
ping Web Audio API a chaque nouvelle notification), (b) toast (une
carte animee glisse depuis le coin superieur droit), (c) push native
(notification systeme lorsque le tableau de bord est installe comme
PWA et que la permission de notification est accordee).

Le compteur de la pastille cloche s'incremente a chaque element non lu ;
cliquer dessus ouvre la page Notifications ou chaque evenement des
trente derniers jours est liste.

\section{PWA patient}

\subsection{Pourquoi une application monopage separee}
Nous avons envisage de faire de l'experience patient une sous-route de
l'application monopage du tableau de bord. Deux arguments nous ont
pousse vers une application separee : le modele cognitif du patient
est completement different de celui du proprietaire (un patient veut
voir son prochain rendez-vous, sa position dans la file, ses fichiers ;
il ne veut pas voir dix-huit elements de barre laterale) ; et
l'installation mobile est une bete differente (les PWA ont besoin de
leur propre manifeste, leur propre service worker, leur propre icone
d'ecran d'accueil).

La PWA reside sur \texttt{/pwa/} sur le meme domaine. Elle a son propre
\texttt{index.html}, son propre \texttt{manifest.json}, son propre
\texttt{sw.js}, sa propre couleur de theme (violet) et son propre
numero de version (actuellement \texttt{nova-pwa-v12}, que nous
incrementons a chaque changement cassant pour que le service worker
recupere).

\subsection{Detection de role}
Un utilisateur se connectant a la PWA peut etre l'une de trois choses :
un patient, un medecin (membre d'equipe avec un \texttt{auth\_uid}), ou
un proprietaire dont le seul appareil est son telephone. Nous detectons
le role au demarrage en interrogeant trois tables en parallele
(\texttt{patient\_users}, \texttt{users}, \texttt{businesses}) et en
stockant le role actif dans \texttt{localStorage}.

\subsection{Onboarding par OTP telephone}
Les patients se connectent par lien magique, mais pour debloquer les
fonctionnalites cote patient, nous avons besoin de leur numero de
telephone --- et nous avons besoin de savoir qu'il est vraiment a eux.
Nous envoyons un code a six chiffres via WhatsApp ; le patient le tape
dans la PWA ; la fonction de verification hache le code (SHA-256 avec
le telephone comme sel), compare au hache stocke, decremente un
compteur d'essai en cas d'echec, et en cas de succes appelle une RPC
Postgres \texttt{pwa\_claim\_phone} qui declenche un trigger
\texttt{link\_clients\_to\_patient\_on\_verify} au niveau base. Ce
trigger est la magie qui connecte l'identite PWA du patient a toutes
les lignes \texttt{clients} eparpillees dans les entreprises qu'il a
visitees.

\subsection{Ecran d'accueil patient}
Apres verification du telephone, le patient atterrit sur une carte hero
montrant la salutation (bonjour / bon apres-midi / bonsoir), le
prochain rendez-vous a venir, une liste de rendez-vous a venir
en-dessous, et des boutons d'action rapide (Trouver une entreprise,
Fichiers). La navigation inferieure a quatre onglets : Accueil,
Decouvrir, Fichiers, Moi.

\subsection{Journee du medecin dans la PWA}
Pour les praticiens, la PWA montre une page d'accueil differente :
les rendez-vous du jour tries par heure, avec actions rapides
\emph{Termine} et \emph{Absent} sur chaque ligne, et un grand bouton
violet \og Scanner le QR du patient \fg{} en haut (lorsque l'entreprise
n'est pas un laboratoire) qui ouvre la camera.

\section{Revue du sprint et tests}
Nous avons teste le Sprint 2 avec une journee simulee complete dans
une vraie clinique. Un patient a ecrit depuis l'exterieur, Nova a
reserve un rendez-vous, le proprietaire a vu la reservation apparaitre
en direct dans le calendrier avec un son, le patient est arrive et a
ete marque \emph{Arrive}, la consultation s'est deroulee et a ete
marquee \emph{Termine}, et la duree reelle s'est affichee sur la ligne
du rendez-vous. En parallele, un partenaire laboratoire a teste l'OCR
de prescriptions avec une pile de quinze prescriptions ; treize ont ete
reconnues correctement au premier passage.

La retrospective a fait remonter une intuition majeure : les
proprietaires de telephones essayant de se connecter a l'URL du
tableau de bord sur un telephone obtenaient une mise en page bureau
non adaptee. Nous avons ajoute le parametre d'echappement
\texttt{forceWeb=1} et la logique de redirection automatique vers la
PWA dans l'iteration suivante.

\section{Conclusion}
A la fin du Sprint 2, chaque ecran que le proprietaire et le patient
touchent au quotidien etait en place. Le sprint final ajoute la couche
d'intelligence qui distingue Nova~AI d'un simple outil de prise de
rendez-vous : remplissage de creneaux, apprentissage des durees,
securite des medicaments, profils vehicules, surprises d'anniversaire,
check-in QR, et file d'attente en temps reel.
